% ---
% titre: Livrets (2-16)
% theme: NO
% chapitre: Nombres naturels et décimaux
% sous_chapitre: Calcul réfléchi et mental
% niveau: [9e]
% type: EX
% tags: [c-livrets, c-multiplication, f-individuel,f-maison,f-station,p-entrainement,p-remediation]
% difficulte: 1
% source: latex
% ---

\header{Livrets de 2, 3 et 4}
 
Complète chaque multiplication en trouvant le résultat manquant.
 
\vfill
\def\x{1.3cm}
\def\arraystretch{2.7}
\begin{tabularx}{\linewidth}{X r c l c l @{\hspace{15pt}}|@{\hspace{15pt}} X r c l c l @{\hspace{15pt}}|@{\hspace{15pt}} X r c l c l}
1) & 2 & $\times$ & 3 & = & \sol{\bl{\x}}{6} & 1) & 3 & $\times$ & 10 & = & \sol{\bl{\x}}{30} & 1) & 4 & $\times$ & 12 & = & \sol{\bl{\x}}{48} \\
2) & 2 & $\times$ & 2 & = & \sol{\bl{\x}}{4} & 2) & 3 & $\times$ & 7 & = & \sol{\bl{\x}}{21} & 2) & 4 & $\times$ & 4 & = & \sol{\bl{\x}}{16} \\
3) & 2 & $\times$ & 14 & = & \sol{\bl{\x}}{28} & 3) & 3 & $\times$ & 6 & = & \sol{\bl{\x}}{18} & 3) & 4 & $\times$ & 6 & = & \sol{\bl{\x}}{24} \\
4) & 2 & $\times$ & 8 & = & \sol{\bl{\x}}{16} & 4) & 3 & $\times$ & 14 & = & \sol{\bl{\x}}{42} & 4) & 4 & $\times$ & 11 & = & \sol{\bl{\x}}{44} \\
5) & 2 & $\times$ & 6 & = & \sol{\bl{\x}}{12} & 5) & 3 & $\times$ & 2 & = & \sol{\bl{\x}}{6} & 5) & 4 & $\times$ & 9 & = & \sol{\bl{\x}}{36} \\
6) & 2 & $\times$ & 9 & = & \sol{\bl{\x}}{18} & 6) & 3 & $\times$ & 4 & = & \sol{\bl{\x}}{12} & 6) & 4 & $\times$ & 3 & = & \sol{\bl{\x}}{12} \\
7) & 2 & $\times$ & 5 & = & \sol{\bl{\x}}{10} & 7) & 3 & $\times$ & 12 & = & \sol{\bl{\x}}{36} & 7) & 4 & $\times$ & 10 & = & \sol{\bl{\x}}{40} \\
8) & 2 & $\times$ & 13 & = & \sol{\bl{\x}}{26} & 8) & 3 & $\times$ & 8 & = & \sol{\bl{\x}}{24} & 8) & 4 & $\times$ & 15 & = & \sol{\bl{\x}}{60} \\
9) & 2 & $\times$ & 10 & = & \sol{\bl{\x}}{20} & 9) & 3 & $\times$ & 11 & = & \sol{\bl{\x}}{33} & 9) & 4 & $\times$ & 7 & = & \sol{\bl{\x}}{28} \\
10) & 2 & $\times$ & 15 & = & \sol{\bl{\x}}{30} & 10) & 3 & $\times$ & 3 & = & \sol{\bl{\x}}{9} & 10) & 4 & $\times$ & 8 & = & \sol{\bl{\x}}{32} \\
11) & 2 & $\times$ & 7 & = & \sol{\bl{\x}}{14} & 11) & 3 & $\times$ & 5 & = & \sol{\bl{\x}}{15} & 11) & 4 & $\times$ & 14 & = & \sol{\bl{\x}}{56} \\
12) & 2 & $\times$ & 11 & = & \sol{\bl{\x}}{22} & 12) & 3 & $\times$ & 15 & = & \sol{\bl{\x}}{45} & 12) & 4 & $\times$ & 13 & = & \sol{\bl{\x}}{52} \\
13) & 2 & $\times$ & 12 & = & \sol{\bl{\x}}{24} & 13) & 3 & $\times$ & 9 & = & \sol{\bl{\x}}{27} & 13) & 4 & $\times$ & 5 & = & \sol{\bl{\x}}{20} \\
14) & 2 & $\times$ & 4 & = & \sol{\bl{\x}}{8} & 14) & 3 & $\times$ & 13 & = & \sol{\bl{\x}}{39} & 14) & 4 & $\times$ & 2 & = & \sol{\bl{\x}}{8} \\
\end{tabularx}
 
\newpage
\header{Livrets de 5, 6 et 7}
 
Complète chaque multiplication en trouvant le résultat manquant.
 
\vfill
\begin{tabularx}{1\linewidth}{X r c l c l @{\hspace{15pt}}|@{\hspace{15pt}} X r c l c l @{\hspace{15pt}}|@{\hspace{15pt}} X r c l c l}
1) & 5 & $\times$ & 5 & = & \sol{\bl{\x}}{25} & 1) & 6 & $\times$ & 6 & = & \sol{\bl{\x}}{36} & 1) & 7 & $\times$ & 7 & = & \sol{\bl{\x}}{49} \\
2) & 5 & $\times$ & 10 & = & \sol{\bl{\x}}{50} & 2) & 6 & $\times$ & 7 & = & \sol{\bl{\x}}{42} & 2) & 7 & $\times$ & 4 & = & \sol{\bl{\x}}{28} \\
3) & 5 & $\times$ & 13 & = & \sol{\bl{\x}}{65} & 3) & 6 & $\times$ & 5 & = & \sol{\bl{\x}}{30} & 3) & 7 & $\times$ & 13 & = & \sol{\bl{\x}}{91} \\
4) & 5 & $\times$ & 9 & = & \sol{\bl{\x}}{45} & 4) & 6 & $\times$ & 3 & = & \sol{\bl{\x}}{18} & 4) & 7 & $\times$ & 2 & = & \sol{\bl{\x}}{14} \\
5) & 5 & $\times$ & 8 & = & \sol{\bl{\x}}{40} & 5) & 6 & $\times$ & 11 & = & \sol{\bl{\x}}{66} & 5) & 7 & $\times$ & 10 & = & \sol{\bl{\x}}{70} \\
6) & 5 & $\times$ & 12 & = & \sol{\bl{\x}}{60} & 6) & 6 & $\times$ & 4 & = & \sol{\bl{\x}}{24} & 6) & 7 & $\times$ & 15 & = & \sol{\bl{\x}}{105} \\
7) & 5 & $\times$ & 4 & = & \sol{\bl{\x}}{20} & 7) & 6 & $\times$ & 8 & = & \sol{\bl{\x}}{48} & 7) & 7 & $\times$ & 3 & = & \sol{\bl{\x}}{21} \\
8) & 5 & $\times$ & 14 & = & \sol{\bl{\x}}{70} & 8) & 6 & $\times$ & 15 & = & \sol{\bl{\x}}{90} & 8) & 7 & $\times$ & 6 & = & \sol{\bl{\x}}{42} \\
9) & 5 & $\times$ & 6 & = & \sol{\bl{\x}}{30} & 9) & 6 & $\times$ & 14 & = & \sol{\bl{\x}}{84} & 9) & 7 & $\times$ & 9 & = & \sol{\bl{\x}}{63} \\
10) & 5 & $\times$ & 2 & = & \sol{\bl{\x}}{10} & 10) & 6 & $\times$ & 12 & = & \sol{\bl{\x}}{72} & 10) & 7 & $\times$ & 8 & = & \sol{\bl{\x}}{56} \\
11) & 5 & $\times$ & 3 & = & \sol{\bl{\x}}{15} & 11) & 6 & $\times$ & 10 & = & \sol{\bl{\x}}{60} & 11) & 7 & $\times$ & 12 & = & \sol{\bl{\x}}{84} \\
12) & 5 & $\times$ & 15 & = & \sol{\bl{\x}}{75} & 12) & 6 & $\times$ & 2 & = & \sol{\bl{\x}}{12} & 12) & 7 & $\times$ & 14 & = & \sol{\bl{\x}}{98} \\
13) & 5 & $\times$ & 7 & = & \sol{\bl{\x}}{35} & 13) & 6 & $\times$ & 9 & = & \sol{\bl{\x}}{54} & 13) & 7 & $\times$ & 11 & = & \sol{\bl{\x}}{77} \\
14) & 5 & $\times$ & 11 & = & \sol{\bl{\x}}{55} & 14) & 6 & $\times$ & 13 & = & \sol{\bl{\x}}{78} & 14) & 7 & $\times$ & 5 & = & \sol{\bl{\x}}{35} \\
\end{tabularx}
 
\newpage
\header{Livrets de 8, 9 et 10}
 
Complète chaque multiplication en trouvant le résultat manquant.
 
\vfill
\begin{tabularx}{1\linewidth}{X r c l c l @{\hspace{15pt}}|@{\hspace{15pt}} X r c l c l @{\hspace{15pt}}|@{\hspace{15pt}} X r c l c l}
1) & 8 & $\times$ & 5 & = & \sol{\bl{\x}}{40} & 1) & 9 & $\times$ & 8 & = & \sol{\bl{\x}}{72} & 1) & 10 & $\times$ & 13 & = & \sol{\bl{\x}}{130} \\
2) & 8 & $\times$ & 10 & = & \sol{\bl{\x}}{80} & 2) & 9 & $\times$ & 11 & = & \sol{\bl{\x}}{99} & 2) & 10 & $\times$ & 3 & = & \sol{\bl{\x}}{30} \\
3) & 8 & $\times$ & 7 & = & \sol{\bl{\x}}{56} & 3) & 9 & $\times$ & 7 & = & \sol{\bl{\x}}{63} & 3) & 10 & $\times$ & 4 & = & \sol{\bl{\x}}{40} \\
4) & 8 & $\times$ & 9 & = & \sol{\bl{\x}}{72} & 4) & 9 & $\times$ & 4 & = & \sol{\bl{\x}}{36} & 4) & 10 & $\times$ & 12 & = & \sol{\bl{\x}}{120} \\
5) & 8 & $\times$ & 15 & = & \sol{\bl{\x}}{120} & 5) & 9 & $\times$ & 14 & = & \sol{\bl{\x}}{126} & 5) & 10 & $\times$ & 2 & = & \sol{\bl{\x}}{20} \\
6) & 8 & $\times$ & 2 & = & \sol{\bl{\x}}{16} & 6) & 9 & $\times$ & 15 & = & \sol{\bl{\x}}{135} & 6) & 10 & $\times$ & 10 & = & \sol{\bl{\x}}{100} \\
7) & 8 & $\times$ & 14 & = & \sol{\bl{\x}}{112} & 7) & 9 & $\times$ & 12 & = & \sol{\bl{\x}}{108} & 7) & 10 & $\times$ & 7 & = & \sol{\bl{\x}}{70} \\
8) & 8 & $\times$ & 11 & = & \sol{\bl{\x}}{88} & 8) & 9 & $\times$ & 3 & = & \sol{\bl{\x}}{27} & 8) & 10 & $\times$ & 6 & = & \sol{\bl{\x}}{60} \\
9) & 8 & $\times$ & 6 & = & \sol{\bl{\x}}{48} & 9) & 9 & $\times$ & 2 & = & \sol{\bl{\x}}{18} & 9) & 10 & $\times$ & 9 & = & \sol{\bl{\x}}{90} \\
10) & 8 & $\times$ & 8 & = & \sol{\bl{\x}}{64} & 10) & 9 & $\times$ & 10 & = & \sol{\bl{\x}}{90} & 10) & 10 & $\times$ & 15 & = & \sol{\bl{\x}}{150} \\
11) & 8 & $\times$ & 12 & = & \sol{\bl{\x}}{96} & 11) & 9 & $\times$ & 13 & = & \sol{\bl{\x}}{117} & 11) & 10 & $\times$ & 5 & = & \sol{\bl{\x}}{50} \\
12) & 8 & $\times$ & 3 & = & \sol{\bl{\x}}{24} & 12) & 9 & $\times$ & 9 & = & \sol{\bl{\x}}{81} & 12) & 10 & $\times$ & 11 & = & \sol{\bl{\x}}{110} \\
13) & 8 & $\times$ & 13 & = & \sol{\bl{\x}}{104} & 13) & 9 & $\times$ & 5 & = & \sol{\bl{\x}}{45} & 13) & 10 & $\times$ & 14 & = & \sol{\bl{\x}}{140} \\
14) & 8 & $\times$ & 4 & = & \sol{\bl{\x}}{32} & 14) & 9 & $\times$ & 6 & = & \sol{\bl{\x}}{54} & 14) & 10 & $\times$ & 8 & = & \sol{\bl{\x}}{80} \\
\end{tabularx}
 
\newpage
\header{Livrets de 11, 12 et 13}
 
Complète chaque multiplication en trouvant le résultat manquant.
 
\vfill
\begin{tabularx}{1\linewidth}{X r c l c l @{\hspace{15pt}}|@{\hspace{15pt}} X r c l c l @{\hspace{15pt}}|@{\hspace{15pt}} X r c l c l}
1) & 11 & $\times$ & 12 & = & \sol{\bl{\x}}{132} & 1) & 12 & $\times$ & 5 & = & \sol{\bl{\x}}{60} & 1) & 13 & $\times$ & 4 & = & \sol{\bl{\x}}{52} \\
2) & 11 & $\times$ & 2 & = & \sol{\bl{\x}}{22} & 2) & 12 & $\times$ & 14 & = & \sol{\bl{\x}}{168} & 2) & 13 & $\times$ & 10 & = & \sol{\bl{\x}}{130} \\
3) & 11 & $\times$ & 4 & = & \sol{\bl{\x}}{44} & 3) & 12 & $\times$ & 2 & = & \sol{\bl{\x}}{24} & 3) & 13 & $\times$ & 8 & = & \sol{\bl{\x}}{104} \\
4) & 11 & $\times$ & 11 & = & \sol{\bl{\x}}{121} & 4) & 12 & $\times$ & 4 & = & \sol{\bl{\x}}{48} & 4) & 13 & $\times$ & 14 & = & \sol{\bl{\x}}{182} \\
5) & 11 & $\times$ & 13 & = & \sol{\bl{\x}}{143} & 5) & 12 & $\times$ & 3 & = & \sol{\bl{\x}}{36} & 5) & 13 & $\times$ & 11 & = & \sol{\bl{\x}}{143} \\
6) & 11 & $\times$ & 3 & = & \sol{\bl{\x}}{33} & 6) & 12 & $\times$ & 7 & = & \sol{\bl{\x}}{84} & 6) & 13 & $\times$ & 5 & = & \sol{\bl{\x}}{65} \\
7) & 11 & $\times$ & 6 & = & \sol{\bl{\x}}{66} & 7) & 12 & $\times$ & 10 & = & \sol{\bl{\x}}{120} & 7) & 13 & $\times$ & 3 & = & \sol{\bl{\x}}{39} \\
8) & 11 & $\times$ & 10 & = & \sol{\bl{\x}}{110} & 8) & 12 & $\times$ & 8 & = & \sol{\bl{\x}}{96} & 8) & 13 & $\times$ & 7 & = & \sol{\bl{\x}}{91} \\
9) & 11 & $\times$ & 8 & = & \sol{\bl{\x}}{88} & 9) & 12 & $\times$ & 15 & = & \sol{\bl{\x}}{180} & 9) & 13 & $\times$ & 9 & = & \sol{\bl{\x}}{117} \\
10) & 11 & $\times$ & 14 & = & \sol{\bl{\x}}{154} & 10) & 12 & $\times$ & 6 & = & \sol{\bl{\x}}{72} & 10) & 13 & $\times$ & 15 & = & \sol{\bl{\x}}{195} \\
11) & 11 & $\times$ & 9 & = & \sol{\bl{\x}}{99} & 11) & 12 & $\times$ & 13 & = & \sol{\bl{\x}}{156} & 11) & 13 & $\times$ & 6 & = & \sol{\bl{\x}}{78} \\
12) & 11 & $\times$ & 7 & = & \sol{\bl{\x}}{77} & 12) & 12 & $\times$ & 11 & = & \sol{\bl{\x}}{132} & 12) & 13 & $\times$ & 13 & = & \sol{\bl{\x}}{169} \\
13) & 11 & $\times$ & 5 & = & \sol{\bl{\x}}{55} & 13) & 12 & $\times$ & 12 & = & \sol{\bl{\x}}{144} & 13) & 13 & $\times$ & 12 & = & \sol{\bl{\x}}{156} \\
14) & 11 & $\times$ & 15 & = & \sol{\bl{\x}}{165} & 14) & 12 & $\times$ & 9 & = & \sol{\bl{\x}}{108} & 14) & 13 & $\times$ & 2 & = & \sol{\bl{\x}}{26} \\
\end{tabularx}
 
\newpage
\header{Livrets de 14, 15 et 16}
 
Complète chaque multiplication en trouvant le résultat manquant.
 
\vfill
\begin{tabularx}{1\linewidth}{X r c l c l @{\hspace{15pt}}|@{\hspace{15pt}} X r c l c l @{\hspace{15pt}}|@{\hspace{15pt}} X r c l c l}
1) & 14 & $\times$ & 13 & = & \sol{\bl{\x}}{182} & 1) & 15 & $\times$ & 14 & = & \sol{\bl{\x}}{210} & 1) & 16 & $\times$ & 4 & = & \sol{\bl{\x}}{64} \\
2) & 14 & $\times$ & 10 & = & \sol{\bl{\x}}{140} & 2) & 15 & $\times$ & 2 & = & \sol{\bl{\x}}{30} & 2) & 16 & $\times$ & 8 & = & \sol{\bl{\x}}{128} \\
3) & 14 & $\times$ & 14 & = & \sol{\bl{\x}}{196} & 3) & 15 & $\times$ & 11 & = & \sol{\bl{\x}}{165} & 3) & 16 & $\times$ & 12 & = & \sol{\bl{\x}}{192} \\
4) & 14 & $\times$ & 9 & = & \sol{\bl{\x}}{126} & 4) & 15 & $\times$ & 9 & = & \sol{\bl{\x}}{135} & 4) & 16 & $\times$ & 14 & = & \sol{\bl{\x}}{224} \\
5) & 14 & $\times$ & 12 & = & \sol{\bl{\x}}{168} & 5) & 15 & $\times$ & 4 & = & \sol{\bl{\x}}{60} & 5) & 16 & $\times$ & 5 & = & \sol{\bl{\x}}{80} \\
6) & 14 & $\times$ & 8 & = & \sol{\bl{\x}}{112} & 6) & 15 & $\times$ & 8 & = & \sol{\bl{\x}}{120} & 6) & 16 & $\times$ & 9 & = & \sol{\bl{\x}}{144} \\
7) & 14 & $\times$ & 6 & = & \sol{\bl{\x}}{84} & 7) & 15 & $\times$ & 12 & = & \sol{\bl{\x}}{180} & 7) & 16 & $\times$ & 7 & = & \sol{\bl{\x}}{112} \\
8) & 14 & $\times$ & 4 & = & \sol{\bl{\x}}{56} & 8) & 15 & $\times$ & 15 & = & \sol{\bl{\x}}{225} & 8) & 16 & $\times$ & 3 & = & \sol{\bl{\x}}{48} \\
9) & 14 & $\times$ & 2 & = & \sol{\bl{\x}}{28} & 9) & 15 & $\times$ & 7 & = & \sol{\bl{\x}}{105} & 9) & 16 & $\times$ & 2 & = & \sol{\bl{\x}}{32} \\
10) & 14 & $\times$ & 15 & = & \sol{\bl{\x}}{210} & 10) & 15 & $\times$ & 13 & = & \sol{\bl{\x}}{195} & 10) & 16 & $\times$ & 6 & = & \sol{\bl{\x}}{96} \\
11) & 14 & $\times$ & 7 & = & \sol{\bl{\x}}{98} & 11) & 15 & $\times$ & 10 & = & \sol{\bl{\x}}{150} & 11) & 16 & $\times$ & 10 & = & \sol{\bl{\x}}{160} \\
12) & 14 & $\times$ & 3 & = & \sol{\bl{\x}}{42} & 12) & 15 & $\times$ & 3 & = & \sol{\bl{\x}}{45} & 12) & 16 & $\times$ & 11 & = & \sol{\bl{\x}}{176} \\
13) & 14 & $\times$ & 11 & = & \sol{\bl{\x}}{154} & 13) & 15 & $\times$ & 5 & = & \sol{\bl{\x}}{75} & 13) & 16 & $\times$ & 13 & = & \sol{\bl{\x}}{208} \\
14) & 14 & $\times$ & 5 & = & \sol{\bl{\x}}{70} & 14) & 15 & $\times$ & 6 & = & \sol{\bl{\x}}{90} & 14) & 16 & $\times$ & 15 & = & \sol{\bl{\x}}{240} \\
\end{tabularx}
 
\newpage
\header{Mélange de tout (livrets 2 à 16)}
 
Complète chaque multiplication en trouvant le résultat manquant.
 
\vfill
\begin{tabularx}{1\linewidth}{X r c l c l @{\hspace{15pt}}|@{\hspace{15pt}} X r c l c l @{\hspace{15pt}}|@{\hspace{15pt}} X r c l c l}
1) & 14 & $\times$ & 16 & = & \sol{\bl{\x}}{224} & 1) & 10 & $\times$ & 11 & = & \sol{\bl{\x}}{110} & 1) & 9 & $\times$ & 10 & = & \sol{\bl{\x}}{90} \\
2) & 2 & $\times$ & 14 & = & \sol{\bl{\x}}{28} & 2) & 8 & $\times$ & 16 & = & \sol{\bl{\x}}{128} & 2) & 7 & $\times$ & 5 & = & \sol{\bl{\x}}{35} \\
3) & 14 & $\times$ & 5 & = & \sol{\bl{\x}}{70} & 3) & 5 & $\times$ & 4 & = & \sol{\bl{\x}}{20} & 3) & 13 & $\times$ & 9 & = & \sol{\bl{\x}}{117} \\
4) & 15 & $\times$ & 3 & = & \sol{\bl{\x}}{45} & 4) & 4 & $\times$ & 6 & = & \sol{\bl{\x}}{24} & 4) & 9 & $\times$ & 10 & = & \sol{\bl{\x}}{90} \\
5) & 10 & $\times$ & 16 & = & \sol{\bl{\x}}{160} & 5) & 9 & $\times$ & 4 & = & \sol{\bl{\x}}{36} & 5) & 9 & $\times$ & 8 & = & \sol{\bl{\x}}{72} \\
6) & 8 & $\times$ & 14 & = & \sol{\bl{\x}}{112} & 6) & 12 & $\times$ & 11 & = & \sol{\bl{\x}}{132} & 6) & 2 & $\times$ & 13 & = & \sol{\bl{\x}}{26} \\
7) & 4 & $\times$ & 8 & = & \sol{\bl{\x}}{32} & 7) & 11 & $\times$ & 3 & = & \sol{\bl{\x}}{33} & 7) & 7 & $\times$ & 3 & = & \sol{\bl{\x}}{21} \\
8) & 13 & $\times$ & 7 & = & \sol{\bl{\x}}{91} & 8) & 13 & $\times$ & 8 & = & \sol{\bl{\x}}{104} & 8) & 9 & $\times$ & 10 & = & \sol{\bl{\x}}{90} \\
9) & 11 & $\times$ & 12 & = & \sol{\bl{\x}}{132} & 9) & 10 & $\times$ & 2 & = & \sol{\bl{\x}}{20} & 9) & 9 & $\times$ & 10 & = & \sol{\bl{\x}}{90} \\
10) & 14 & $\times$ & 8 & = & \sol{\bl{\x}}{112} & 10) & 5 & $\times$ & 5 & = & \sol{\bl{\x}}{25} & 10) & 3 & $\times$ & 16 & = & \sol{\bl{\x}}{48} \\
11) & 14 & $\times$ & 11 & = & \sol{\bl{\x}}{154} & 11) & 14 & $\times$ & 2 & = & \sol{\bl{\x}}{28} & 11) & 6 & $\times$ & 15 & = & \sol{\bl{\x}}{90} \\
12) & 5 & $\times$ & 11 & = & \sol{\bl{\x}}{55} & 12) & 8 & $\times$ & 2 & = & \sol{\bl{\x}}{16} & 12) & 7 & $\times$ & 11 & = & \sol{\bl{\x}}{77} \\
13) & 16 & $\times$ & 2 & = & \sol{\bl{\x}}{32} & 13) & 11 & $\times$ & 10 & = & \sol{\bl{\x}}{110} & 13) & 16 & $\times$ & 10 & = & \sol{\bl{\x}}{160} \\
14) & 2 & $\times$ & 13 & = & \sol{\bl{\x}}{26} & 14) & 12 & $\times$ & 4 & = & \sol{\bl{\x}}{48} & 14) & 13 & $\times$ & 12 & = & \sol{\bl{\x}}{156} \\
\end{tabularx}
\def\arraystretch{1}